\documentclass[11pt]{article}
\usepackage[margin=1in]{geometry}
\usepackage{amsmath,amssymb}
\usepackage{enumitem}
\setlength{\parindent}{0pt}
\setlength{\parskip}{6pt}
\begin{document}
\section*{STAT 412 HW07 --- Question 12}

\subsection*{Problem}

\begin{quote}
It is claimed that automobiles are driven on average more than 18,000 kilometers per year. To test this claim, 110 randomly selected automobile owners are asked to keep a record of the kilometers they travel. Would you agree with this claim if the random sample showed an average of 18,330 kilometers and a standard deviation of 3700 kilometers? Use a P-value in your conclusion.
\end{quote}

\subsection*{Solution}

The claim is ``more than 18,000,'' so the alternative is right-tailed:
\[
H_0:\mu=18000,\qquad H_a:\mu>18000.
\]
This corresponds to choice D.

The population standard deviation is not given, so use a one-sample $t$ statistic:
\[
t=\frac{\bar{x}-\mu_0}{s/\sqrt{n}}
  =\frac{18330-18000}{3700/\sqrt{110}}.
\]

Compute the denominator:
\[
\frac{3700}{\sqrt{110}}\approx 352.78.
\]
Then
\[
t=\frac{330}{352.78}\approx 0.93542.
\]

Rounded to two decimals, the test statistic is
\[
t=0.94.
\]

With $df=109$, the right-tail P-value is about
\[
P(T_{109}>0.93542)\approx 0.176.
\]

Since the P-value is large, do not reject $H_0$. The sample mean is above 18,000, but not by enough to statistically support the ``more than 18,000'' claim.

\subsection*{Final Answer}

\fbox{\begin{minipage}{0.94\linewidth}
Choice D: $H_0:\mu=18000$; $H_a:\mu>18000$. Test statistic $t=0.94$; P-value $\approx0.176$. Do not reject $H_0$; there is not enough evidence to agree with the claim.
\end{minipage}}

\end{document}
